\chapter{Eating Disorders and Body Image}
\index{Eating Disorders and Body Image}

\section*{Definition and Overview}
Body image is how a person sees, thinks, and feels about their body, and it sits at the centre of many eating disorders. A negative body image involves persistent dissatisfaction with appearance, distorted perceptions of size or shape, and the belief that appearance defines self-worth. Body dissatisfaction is a powerful driver of dieting and disordered eating, and it is one of the strongest risk factors for the development of eating disorders. This chapter explores the link between body image and eating disorders, who is affected, and what can be done to promote a healthier relationship with body and food.

\section*{Epidemiology}
Body dissatisfaction is widespread. A large proportion of adolescent girls and a substantial number of boys report wanting to change their bodies, and body dissatisfaction is increasingly common in adults as well. It affects people across all body sizes and is linked to dieting, disordered eating, and mental health problems. Social media use and exposure to idealised images are associated with poorer body image, particularly in young people.

\section*{Aetiology and Pathophysiology}
Body image is shaped by a complex mix of influences.
\begin{itemize}
    \item \textbf{Sociocultural pressures:} media, advertising, and social media present narrow and often unattainable body ideals
    \item \textbf{Social influences:} comments from family and peers, weight teasing, and comparison with others
    \item \textbf{Psychological factors:} perfectionism, low self-esteem, and the tendency to compare oneself with others
    \item \textbf{Internalisation:} adopting society's thin ideal as one's own standard
    \item \textbf{Body comparison:} constant comparison with others, online and in person, worsens dissatisfaction
    \item \textbf{The pathway to eating disorders:} body dissatisfaction leads to dieting, and dieting can spiral into disordered eating and a full eating disorder in vulnerable people
    \item \textbf{The effects of disordered eating:} starvation and weight changes further distort body image, creating a vicious cycle
\end{itemize}

\section*{Clinical Features}
The signs of a negative body image and its progression towards an eating disorder include the following.
\begin{itemize}
    \item Constant criticism of one's body and preoccupation with weight, shape, and appearance
    \item Checking behaviour: repeatedly weighing, measuring, or looking in the mirror
    \item Avoidance: avoiding mirrors, photos, swimming, and social situations
    \item Comparing one's body with others, including images on social media
    \item Believing that being thinner would solve problems or increase worth
    \item Dieting, skipping meals, or rigid food rules
    \item Excessive exercise to control weight
    \item In an eating disorder: the full features of anorexia, bulimia, or binge eating disorder
\end{itemize}

\section*{Differential Diagnosis}
Body image concerns can be part of several conditions.
\begin{itemize}
    \item Body dysmorphic disorder, in which perceived flaws in appearance dominate thinking
    \item Eating disorders, in which body image disturbance drives disordered eating
    \item Muscle dysmorphia, a preoccupation with being insufficiently muscular
    \item Depression and anxiety, which can include negative self-perception
    \item Normal developmental concerns about appearance, particularly in adolescence
\end{itemize}

\section*{When to Seek Medical Care}
Professional help should be sought when body image concerns cause distress, interfere with daily life, or are accompanied by dieting, weight changes, binge eating, purging, or low mood. Family members should take signs of an eating disorder seriously and encourage early assessment.

\textbf{Contact your local emergency services for:}
\begin{itemize}
    \item Fainting, collapse, or severe weakness in the context of restricted eating
    \item Chest pain or palpitations
    \item Self-harm or thoughts of suicide
    \item Sudden severe abdominal pain
\end{itemize}

\section*{Diagnosis and Investigations}
Assessment explores the relationship between body image, eating, and psychological wellbeing.
\begin{itemize}
    \item \textbf{History:} body image concerns, eating patterns, weight history, exercise, and mental state
    \item \textbf{Questionnaires:} validated tools measure body dissatisfaction and disordered eating
    \item \textbf{Physical assessment:} weight, body mass index, and general health where an eating disorder is present
    \item \textbf{Mental health assessment:} screening for depression, anxiety, and eating disorder symptoms
\end{itemize}

\section*{Management}
Improving body image is central to treating and preventing eating disorders.
\begin{itemize}
    \item \textbf{Psychological therapy:} cognitive behavioural therapy addresses the thoughts, behaviours, and avoidance that maintain negative body image
    \item \textbf{Media literacy:} learning to recognise and resist unrealistic images and messages
    \item \textbf{Reducing comparison:} limiting social media and unfollowing accounts that trigger comparison
    \item \textbf{Self-compassion:} practising kindness toward oneself rather than self-criticism
    \item \textbf{Focus on function:} valuing what the body can do rather than how it looks
    \item \textbf{Family approaches:} parents modelling positive body attitudes and avoiding weight talk
    \item \textbf{Treatment of eating disorders:} where an eating disorder is present, full treatment as described in the eating disorders chapter
\end{itemize}

\section*{Complications}
\begin{itemize}
    \item Progression to a full eating disorder
    \item Depression, anxiety, and reduced self-esteem
    \item Avoidance of social, school, and work activities
    \item Exercise dependence and overtraining injuries
    \item Reduced quality of life and wellbeing
\end{itemize}

\section*{Prognosis and Outcomes}
Body image can improve with the right support, and most people who receive therapy for body dissatisfaction or eating disorders make meaningful progress. Early intervention before the development of a full eating disorder is especially effective. Recovery involves learning to see oneself with less distortion and more compassion, and it is achievable with time and support.

\section*{Prevention and Self-Care}
\begin{itemize}
    \item Practise media literacy and question unrealistic images
    \item Reduce social media use and curate a positive feed
    \item Avoid weighing yourself frequently and ditch the scales where helpful
    \item Focus on health, energy, and what your body can do
    \item Avoid negative body talk, both your own and others'
    \item Surround yourself with people who do not focus on weight
    \item Seek help early if body concerns are affecting your life
\end{itemize}

\section*{Sources and Further Reading}
The following reputable sources provide further information for healthcare students and patients seeking reliable, current advice about body image and eating disorders.
\begin{itemize}
    \item Butterfly Foundation. \textit{Body image information and support.} https://butterfly.org.au/
    \item National Eating Disorders Collaboration. \textit{Body image resources.} https://nedc.com.au/
    \item ReachOut Australia. \textit{Body image resources for young people.} https://au.reachout.com/
    \item NHS. \textit{Body image: mental health information.} https://www.nhs.uk/mental-health/
    \item Head to Health. \textit{Body image and eating concerns.} https://www.headtohealth.gov.au/
    \item Eating Disorders Victoria. \textit{Resources on body image.} https://www.eatingdisorders.org.au/
\end{itemize}
